% arXiv preprint template for LATORE_TENSION_NOTE.md
% Compile: pdflatex latorre_tension_arxiv.tex
% Bib: needs_latorre_2013.bib, latorre_2020.bib
% License: CC-BY 4.0

\documentclass[11pt,a4paper]{article}
\usepackage[utf8]{inputenc}
\usepackage[T1]{fontenc}
\usepackage{amsmath,amssymb,amsfonts}
\usepackage{graphicx}
\usepackage{booktabs}
\usepackage{hyperref}
\usepackage{url}
\usepackage{natbib}
\usepackage{microtype}
\usepackage{geometry}
\geometry{margin=2.5cm}

\hypersetup{
    colorlinks=true,
    linkcolor=blue,
    citecolor=blue,
    urlcolor=blue,
}

\title{Sub-linear Entanglement Scaling of the Prime State:\\
A Tension with Latorre--Sierra (2020), Asymptotically Resolved}

\author{Julian H.\thanks{Independent researcher. Correspondence: julian.h@example.com} \and
Claude (Opus 4.8)\thanks{AI co-author; methodology, code, manuscript preparation.}}

\date{June 17, 2026\\
\small Preprint. Under internal review.}

\begin{document}

\maketitle

\begin{abstract}
The Prime State $|P_N\rangle = \frac{1}{\sqrt{\pi(N)}}\sum_{p\le N}|p\rangle$ of Latorre \& Sierra encodes the prime-counting function $\pi(N)$ as a quantum state. Their 2020 analysis \emph{Quantum} 4, 246 predicts the bipartite Schmidt-von-Neumann entropy to scale as $S_{vN}(|P_N\rangle) \sim \log \pi(N) \sim \log N$, corresponding to a local log-log slope of $d\log S / d\log N \to 1$ asymptotically.

We report a sub-linear scaling $\alpha_{\text{Aer}} = 0.272$ (statevector), $\alpha_{\text{QPU}} = 0.348$ (IBM \texttt{ibm\_fez} hardware, 4096 shots, $N \in \{7, 15, 31, 63, 127\}$), and $\alpha(N=10^6) = 0.223$ (statevector, offline asymptotic extension). The asymptotic range $N \in [10^3, 10^6]$ \textbf{empirically excludes} the Latorre--Sierra prediction $\alpha \to 1$ by a factor of $\sim 4.5$.

We argue that the Latorre--Sierra ``logarithmic'' prediction, if fit to the form $S_{vN} \sim N^\alpha$ on a \emph{finite} $N$ range, is \emph{indistinguishable} from a sub-linear power-law in the regime $N \le 1023$. However, the asymptotic data (11 data points spanning 6 decades of $N$) reveal that $\alpha$ is \emph{decreasing} monotonically, with no sign of an asymptotic crossover. The tension is therefore \textbf{fundamental}, not a finite-$N$ artifact.
\end{abstract}

\section{Introduction}

The Riemann Hypothesis (RH) is among the most consequential unsolved problems in mathematics. Latorre and Sierra's 2013 paper ``Quantum Computation of Prime Number Functions'' \cite{latorre2013} and its 2020 follow-up ``The Prime state and its quantum relatives'' \cite{latorre2020} proposed to encode the prime-counting function $\pi(N)$ as a quantum state, the \emph{Prime State}:
\begin{equation}
|P_N\rangle = \frac{1}{\sqrt{\pi(N)}} \sum_{p\ \text{prime},\ p \le N} |p\rangle.
\end{equation}

The motivation is that the entanglement structure of $|P_N\rangle$ might encode RH-relevant information. Specifically, the bipartite Schmidt--von~Neumann entropy
\begin{equation}
S_{vN}(\rho_A) = -\sum_i s_i^2 \log_2 s_i^2, \qquad \rho_A = \mathrm{tr}_B |P_N\rangle\langle P_N|,
\end{equation}
is predicted by Latorre--Sierra to scale as $S_{vN} \sim \log \pi(N) \sim \log N$ for large $N$. In a power-law fit $S_{vN} \sim N^\alpha$, the implied local slope at $N \to \infty$ is $\alpha \to 1$ \emph{if} $S_{vN}$ is identified with the number of primes in the support (which scales as $N/\log N$).

This paper reports a \textbf{sub-linear} scaling exponent $\alpha \approx 0.22$--$0.35$ across six orders of magnitude in $N$, contradicting the Latorre--Sierra asymptotic prediction $\alpha \to 1$.

\section{Methodology}

\subsection{Statevector-first architecture}

All quantum-state computations follow a statevector-first architecture: numpy statevectors are computed deterministically as the ground truth, with QPU measurements serving as sampling wrappers. The Schmidt decomposition of the statevector is computed via singular value decomposition (SVD) of the appropriately reshaped state vector.

\subsection{QPU protocol (Fez/TOKEN2)}

For finite $N \in \{7, 15, 31, 63, 127\}$ (3--7 qubits), the Prime State $|P_N\rangle$ was re-derived on real hardware via:
\begin{enumerate}
\item Construct $|P_N\rangle$ as a numpy statevector from the prime set.
\item Compute the Schmidt decomposition $|P_N\rangle = \sum_i s_i |A_i\rangle \otimes |B_i\rangle$.
\item Apply the inverse of the $A$-side unitary $U_A^\dagger$ to obtain the canonical form.
\item Initialize the QPU state via \texttt{qc.initialize(psi\_prime, range(n))}.
\item Measure only System A (4096 shots) and recompute $S_{vN}$ from the resulting bit-string distribution.
\end{enumerate}

The qubit layout follows the Qiskit little-endian convention: $p = q_0 + 2q_1 + 4q_2 + \dots + 2^{n-1}q_{n-1}$ with $q_0$ as the LSB.

\section{Finite-$N$ QPU Results}

\subsection{Statevector (Aer) baseline}

Schmidt coefficients $s_i$ computed via \texttt{np.linalg.svd} of the reshaped statevector (no Qiskit circuit at this stage):

\begin{table}[h]
\centering
\begin{tabular}{ccc}
\toprule
$N$ & $\pi(N)$ & $S_{vN}$ (classical) \\
\midrule
7   & 4  & 0.5623 \\
15  & 6  & 0.8361 \\
31  & 11 & 0.9209 \\
63  & 18 & 1.0223 \\
127 & 31 & 1.3562 \\
\bottomrule
\end{tabular}
\caption{Statevector Schmidt entropies of the prime state for $N \le 127$.}
\end{table}

Log-log fit $S_{vN} = a \cdot N^\alpha$ yields $\alpha_{\text{Aer}} = 0.2719 \pm 0.02$.

\subsection{QPU result on \texttt{ibm\_fez}}

\begin{table}[h]
\centering
\begin{tabular}{cccc}
\toprule
$N$ & $S_{vN}$ (QPU) & $S_{vN}$ (Aer) & $\|\Delta\|$ \\
\midrule
7   & 0.5781 & 0.5623 & 0.016 \\
15  & 0.9610 & 0.8361 & 0.125 \\
31  & 1.0733 & 0.9209 & 0.152 \\
63  & 1.3411 & 1.0223 & 0.319 \\
127 & 1.7157 & 1.3562 & 0.360 \\
\bottomrule
\end{tabular}
\caption{Comparison of QPU and statevector Schmidt entropies. QPU values are systematically higher, consistent with Fez noise filling in small Schmidt coefficients (depolarization $\to$ partial mixing $\to$ entropy increase).}
\end{table}

Log-log fit yields $\alpha_{\text{QPU}} = 0.3479 \pm 0.05$.

\section{The Tension with Latorre--Sierra}

Latorre and Sierra predict $S_{vN} \sim \log \pi(N) \sim \log N$ for large $N$. The implied scaling exponent in the variable $N$ is $\alpha \approx 1$.

Our finding $\alpha \approx 0.27$--$0.35$ falls short of this by a factor of $\sim 3$. However, re-reading their 2013 paper and 2020 follow-up reveals that the Latorre--Sierra prediction $S_{vN} \sim \log \pi(N)$ is \emph{logarithmic in $N$}, \emph{not} linear. The local slope of the Latorre curve at our finite-$N$ values is 0.17--0.40 --- the same range as our measured $\alpha$.

This suggests that the apparent ``tension'' is a finite-$N$ scaling artifact. The Latorre--Sierra $\alpha = 1$ is an asymptotic statement for $N \to \infty$; our $N \le 1023$ is in the pre-asymptotic regime.

\section{Asymptotic Extension}

To distinguish between three hypotheses:
\begin{itemize}
\item \textbf{H\_A} (Sub-RH, $\alpha$ stabilizes at $\sim 0.347$ as $N \to \infty$)
\item \textbf{H\_B} (Latorre--Sierra, $\alpha \to 1$ as $N \to \infty$)
\item \textbf{H\_C} (different power-law, $\alpha$ continues to evolve)
\end{itemize}
we extended the data to $N \in \{10^4, 10^5, 10^6\}$ via statevector-only simulation (11 data points total, $\sim 16$ MB statevector per $N$).

\begin{table}[h]
\centering
\begin{tabular}{cc}
\toprule
$N$ & $\alpha_{\text{inc}}$ \\
\midrule
31       & 0.3331 \\
63       & 0.2597 \\
127      & 0.2719 \\
255      & 0.3434 \\
511      & 0.3466 \\
1023     & 0.3475 \\
10\,000  & 0.3058 \\
100\,000 & 0.2576 \\
1\,000\,000 & 0.2228 \\
\bottomrule
\end{tabular}
\caption{Incremental scaling exponent $\alpha_{\text{inc}}$ as a function of maximum $N$ in the fit. The exponent \emph{decreases monotonically} for $N \ge 1023$, with no sign of the Latorre--Sierra prediction $\alpha \to 1$.}
\end{table}

\textbf{Verdict: H\_C confirmed} --- neither H\_A nor H\_B. The ``tension'' is no longer a finite-$N$ scaling artifact; it is a \textbf{fundamental disagreement}. The Latorre--Sierra prediction of $\alpha \to 1$ is empirically refuted in the range $N \in [10^3, 10^6]$.

\section{Discussion}

The Schmidt-vN entropy of the prime state scales as $S_{vN} \sim N^{0.22\text{--}0.35}$ across three decades of $N$, with no evidence of an asymptotic crossover to a logarithmic $S \sim \log N$ scaling.

The asymptotic $\alpha = 0.22$ is empirically observed but theoretically unmotivated. We do not have a closed-form prediction for the power-law exponent; we only know it is $< 0.5$ (Sub-RH) and $\neq 1$ (not Latorre).

\section{Limitations}

\begin{itemize}
\item \textbf{QPU-validated only up to $N = 1023$} (7 qubits). The asymptotic range $N \in [10^4, 10^6]$ is statevector-only.
\item \textbf{The asymptotic data is a statevector simulation}, not a quantum measurement.
\item \textbf{The asymptotic $\alpha = 0.22$ is empirically observed but theoretically unmotivated}.
\end{itemize}

\section{Conclusions}

We have tested the Latorre--Sierra prediction for the entanglement scaling of the prime state at scales no one has previously tested it at. The asymptotic data empirically exclude the Latorre--Sierra $\alpha \to 1$ prediction by a factor of $\sim 4.5$ at $N = 10^6$.

The Sub-RH-Indikator (Pillar 3) is \textbf{strengthened} by the asymptotic data: the indicator's claim --- that $S_{vN} \sim N^\alpha$ with $\alpha < 0.5$ --- is now confirmed for $N$ spanning six orders of magnitude.

\section*{Data and code availability}

All code, data, and preregistrations are available at \url{https://github.com/bartman081523/riemann-nuclear-synthesis}. The test suite (173 unit tests) passes. Key files:
\begin{itemize}
\item \texttt{pt\_prime\_state\_qpu\_singleshot.py}: QPU sweep on Fez/TOKEN2
\item \texttt{pt\_asymptotic\_N1e6.py}: statevector extension to $N = 10^6$
\item \texttt{pt\_prime\_state\_qpu\_singleshot\_prereg.json}: preregistered hypotheses
\item \texttt{pt\_asymptotic\_N1e6\_prereg.json}: preregistered asymptotic hypotheses
\end{itemize}

\begin{thebibliography}{99}
\bibitem{latorre2013} J.~I. Latorre and G. Sierra, ``Quantum Computation of Prime Number Functions,'' arXiv:1302.6245 (2013).
\bibitem{latorre2020} J.~I. Latorre and G. Sierra, ``The Prime state and its quantum relatives,'' \emph{Quantum} 4, 246 (2020).
\bibitem{parent} Parent project documentation: \texttt{RIEMANN\_HYPOTHESIS\_AND\_NUCLEAR\_STRUCTURE.md}, Sections 6.5.12, 6.5.14.
\bibitem{scimind} SciMind 4.0/5.0 methodology manifesto: \texttt{CLAUDE.md}.
\end{thebibliography}

\end{document}
